% TOP_PROBLEM: {"id": 2304027, "title": "Research Problems in Function Theory — Problem 4.27", "category_id": 8, "category": {"id": 8, "name": "analysis", "display_name": "Analysis", "description": "Limits, continuity, calculus, and function theory.", "slug": "analysis", "order_index": 8, "created_at": "2026-07-31T15:26:25.671Z"}, "status": "open", "problem_number": "AMR-022-4027", "set_id": 13}
% FIRST_POSTED: 2026-10-01
% ENTRY_KIND: statement_only
% UnsolvedMath ID 2304027; source catalogue code AMR-022-4027
% !TeX program = lualatex
% Definition and sources only; this file is not an LLM solution attempt.
\documentclass[11pt,a4paper]{article}
\usepackage[margin=25mm]{geometry}
\usepackage{fontspec}
\setmainfont{lmroman10-regular.otf}[BoldFont=lmroman10-bold.otf,
 ItalicFont=lmroman10-italic.otf,BoldItalicFont=lmroman10-bolditalic.otf]
\usepackage{amsmath,amssymb,mathtools}
\usepackage{unicode-math}
\setmathfont{latinmodern-math.otf}
\AtBeginDocument{\let\setminus\smallsetminus}
\usepackage{xurl,hyperref}
\hypersetup{colorlinks=true,urlcolor=blue}
\setcounter{secnumdepth}{0}
\setlength{\parindent}{0pt}
\setlength{\parskip}{5pt}
\setlength{\emergencystretch}{3em}
\begin{document}
\section*{Research Problems in Function Theory — Problem 4.27}\label{2304027}
{\small UnsolvedMath ID 2304027 · AMR-022-4027 · Analysis · AMR Open Problem Lists}
\subsection{Definitions and mathematical statement}
Let $n\geq1$ and let $p\in\mathbb R[x]$ have degree $n$ and
$n$ distinct rational roots. Is there always a nonzero real
number $t$ for which $p(x)-t$ also has $n$ distinct rational roots?
All roots in both fibers are required to be rational and simple;
one additional rational point in a fiber does not suffice.

The question has the following equivalent arithmetic normalization.
Write $p(x)=c q(x)$, where $c\ne0$ is its leading coefficient and
\[
 q(x)=\prod_{j=1}^{n}(x-a_j),\qquad
 a_1,\ldots,a_n\in\mathbb Q\text{ distinct}.
\]
Must there be $s\in\mathbb Q\setminus\{0\}$ and distinct
$b_1,\ldots,b_n\in\mathbb Q$ such that
\[
 \prod_{j=1}^{n}(x-b_j)=q(x)-s?
\]
Any rational root of $q-s$ forces $s$ itself to be rational,
so this normalization loses no possibility. The desired shift
of the original polynomial is $t=cs$.

Equivalently, the two rational root multisets should have the
same elementary symmetric functions of orders $1,\ldots,n-1$,
but different products. Their sets are automatically disjoint
because the two polynomial levels are distinct. The known
low-degree cases motivate the unrestricted question; the
remaining general target is to establish this property in every
degree $n\geq4$, or exhibit a polynomial for which no such
nonzero shift exists.
\subsection{Short English statement}
Does every polynomial with distinct rational roots have another nonzero level at which all its roots are again distinct and rational?
\subsection{Sources}
\begin{itemize}
\item[S1] ULAM.AI and the UnsolvedMath Contributors, supplied problems.json, record 2304027 (AMR-022-4027), AMR Open Problem Lists. Catalogue identity and source transcription; CC BY 4.0.
\url{https://huggingface.co/datasets/ulamai/UnsolvedMath}.
\item[S2] Hayman - Research Problems in Function Theory (2018), Problem 4.27. Original source indexed by AMR-022-4027.
\url{https://arxiv.org/abs/1809.07200}.
\end{itemize}
\end{document}
